\documentclass[12pt,a4paper]{article}
\usepackage[utf8]{vietnam}
\usepackage{amsmath,amssymb}
\usepackage{booktabs}
\usepackage{longtable}
\usepackage{geometry}
\geometry{top=2cm, bottom=2cm, left=2.5cm, right=2cm}
\usepackage{setspace}
\onehalfspacing
\usepackage{hyperref}
\hypersetup{colorlinks=true,linkcolor=blue,urlcolor=blue}
\usepackage{caption}
\usepackage{float}

\title{\textbf{REPORT ON RECURRENCE RISK ANALYSIS\\FOR THYROID CANCER\\ACCORDING TO ATA 2025 GUIDELINES}}
\author{}
\date{}

\begin{document}
\maketitle
\thispagestyle{empty}
\newpage

\section{Study Sample Characteristics}
Total patients: \textbf{114}. Recurrence events: \textbf{7} (6.1\%).

\subsection{ATA 2025 PTC Risk Stratification}
\begin{table}[H]
\centering
\caption{Patient distribution by ATA 2025 PTC 4-level risk stratification}
\label{tab:ata}
\begin{tabular}{lcc}
\toprule
Risk Stratification (PTC) & N (n=114) & Recurrence rate \\
\midrule
Low                 & 69 (60.5\%)   & 5/69 (7.2\%)   \\
Low-Intermediate    & 2 (1.8\%)     & 1/2 (50.0\%)   \\
Intermediate-High   & 15 (13.2\%)   & 0/15 (0.0\%)   \\
High                & 28 (24.6\%)   & 1/28 (3.6\%)   \\
\bottomrule
\end{tabular}
\caption*{\emph{Note:} The 4-level system (Low / Low-Intermediate / Intermediate-High / High) replaces the older 3-level system per ATA 2025 PTC guidelines. Due to missing histopathological criteria (aggressive histology, margins, vascular invasion), these classifications are based on available variables only.}
\end{table}

\subsection{Statistical Method: Firth Logistic Regression}
Firth's penalized likelihood logistic regression is used instead of standard maximum likelihood estimation. Standard logistic regression fails when a predictor has zero events in one group (complete separation), producing infinite odds ratios. This affects TERT, TP53, and co-mutation variables. Firth's method adds a Jeffreys prior penalty $0.5 \cdot \log|I(\beta)|$ to the log-likelihood, guaranteeing finite estimates for all parameters.

\subsection{Descriptive Statistics of Study Variables}
\begin{longtable}{lcc}
\caption{Descriptive statistics of clinical and genetic variables}
\label{tab:desc}\\
\toprule
Variable & N (n=114) & \% \\
\midrule
\endfirsthead
\multicolumn{3}{c}{\tablename\ \thetable\ -- continued} \\
\toprule
Variable & N (n=114) & \% \\
\midrule
\endhead
\bottomrule
\endfoot
Multifocal ($>$1 nodule)                 & 31  & 27.2 \\
Bilateral involvement               & 32  & 28.1 \\
Tumor size $>$1cm              & 21  & 18.4 \\
Tumor size $>$4cm              & 1   & 0.9 \\
T stage T3+                    & 28  & 24.6 \\
Lymph node metastasis (N+)                 & 59  & 51.8 \\
Lateral neck metastasis (N1b)        & 8   & 7.0 \\
Central LN $\geq$5          & 16  & 14.0 \\
Lateral neck LN (+)                  & 21  & 18.4 \\
BRAF V600E                       & 45  & 39.5 \\
Other BRAF                        & 2   & 1.8 \\
TERT mutation                    & 6   & 5.3 \\
TP53 mutation                    & 2   & 1.8 \\
Any mutation            & 51  & 44.7 \\
Co-mutation ($\geq$2 genes)     & 4   & 3.5 \\
BRAF + TERT                      & 3   & 2.6 \\
BRAF + TP53                      & 0   & 0.0 \\
\end{longtable}

\section{Univariate Analysis Results}
\subsection{Firth Univariate Logistic Regression}
\begin{table}[H]
\centering
\caption{Firth univariate analysis of factors associated with recurrence}
\label{tab:uni}
\begin{tabular}{lcccc}
\toprule
Variable & N & OR (Firth) & 95\% CI (Firth) & p \\
\midrule
Multifocal ($>$1 nodule)           & 114 & 2.170 & 0.495--9.514   & 0.304 \\
Bilateral involvement         & 114 & 2.070 & 0.473--9.060   & 0.334 \\
Tumor size $>$1cm        & 114 & 0.985 & 0.151--6.432   & 0.987 \\
Tumor size $>$4cm        & 114 & --    & -- (excluded)   & --   \\
T stage T3+              & 114 & 0.676 & 0.106--4.320   & 0.679 \\
LN metastasis (N+)           & 114 & 4.414 & 0.707--27.553  & 0.112 \\
Lateral neck metastasis (N1b)  & 114 & 3.092 & 0.409--23.381  & 0.274 \\
Central LN $\geq$5    & 114 & 2.931 & 0.574--14.980  & 0.196 \\
Lateral neck LN (+)            & 114 & 0.985 & 0.151--6.432   & 0.987 \\
BRAF V600E                 & 114 & 1.199 & 0.278--5.176   & 0.808 \\
BRAF (any)              & 114 & 1.110 & 0.257--4.788   & 0.889 \\
TERT mutation              & 114 & 1.041 & 0.043--25.413  & 0.980 \\
TP53 mutation              & 114 & 2.813 & 0.063--124.823 & 0.593 \\
Co-mutation ($\geq$2)   & 114 & 1.533 & 0.054--43.723  & 0.803 \\
BRAF + TERT                & 114 & 1.990 & 0.060--66.025  & 0.700 \\
\bottomrule
\end{tabular}
\end{table}

\emph{Note:} Firth penalized logistic regression resolves complete separation. All variables now have finite OR and 95\% CI. TERT, TP53, and co-mutation estimates are finite but have wide CIs due to low event count (7/114). Tumor size $>$4cm (n=1) and BRAF+TP53 (n=0) excluded due to insufficient positive cases.

\section{Multivariate Analysis Results}

\subsection{Model 1 -- Clinical variables (Firth)}
\begin{table}[H]
\centering
\caption{Multivariate Firth analysis -- Clinical model (N=114, events=7)}
\label{tab:multi1}
\begin{tabular}{lccc}
\toprule
Variable & OR (Firth) & 95\% CI (Firth) & p \\
\midrule
Size $>$1cm     & 0.716 & 0.113--4.526  & 0.723 \\
T3+         & 0.570 & 0.094--3.461  & 0.541 \\
N+          & 4.979 & 0.840--29.503 & \textbf{0.077} \\
\bottomrule
\end{tabular}
\caption*{Hosmer-Lemeshow: $\chi^2=4.07$, df=8, p=0.851}
\end{table}

\subsection{Model 2 -- Genetic variables (Firth)}
\begin{table}[H]
\centering
\caption{Multivariate Firth analysis -- Genetic model (N=114, events=7)}
\label{tab:multi2}
\begin{tabular}{lccc}
\toprule
Variable & OR (Firth) & 95\% CI (Firth) & p \\
\midrule
BRAF V600E & 1.180 & 0.282--4.934 & 0.820 \\
TERT       & 1.040 & 0.045--23.937 & 0.980 \\
\bottomrule
\end{tabular}
\caption*{Hosmer-Lemeshow: $\chi^2=7.86$, df=8, p=0.447}
\end{table}

\subsection{Model 3 -- Combined clinical + genetic (Firth)}
\begin{table}[H]
\centering
\caption{Multivariate Firth analysis -- Combined model (N=114, events=7)}
\label{tab:multi3}
\begin{tabular}{lccc}
\toprule
Variable & OR (Firth) & 95\% CI (Firth) & p \\
\midrule
Size $>$1cm     & 0.710 & 0.119--4.218  & 0.706 \\
T3+         & 0.545 & 0.090--3.297  & 0.508 \\
N+          & 4.650 & 0.863--25.049 & \textbf{0.074} \\
BRAF V600E  & 1.111 & 0.265--4.661  & 0.886 \\
TERT        & 0.978 & 0.043--22.227 & 0.989 \\
\bottomrule
\end{tabular}
\caption*{Hosmer-Lemeshow: $\chi^2=6.37$, df=8, p=0.606}
\end{table}

\subsection{Model 4 -- N+ and co-mutation (Firth)}
\begin{table}[H]
\centering
\caption{Multivariate Firth analysis -- N+ and co-mutation (N=114, events=7)}
\label{tab:multi4}
\begin{tabular}{lccc}
\toprule
Variable & OR (Firth) & 95\% CI (Firth) & p \\
\midrule
N+             & 4.118 & 0.710--23.891 & 0.115 \\
Co-mutation & 2.163 & 0.065--71.702 & 0.666 \\
\bottomrule
\end{tabular}
\caption*{Hosmer-Lemeshow: $\chi^2=4.11$, df=8, p=0.848}
\end{table}

\section{Detailed Analysis of BRAF V600E}
BRAF V600E prevalence is 45/114 (39.5\%). Univariate Firth OR=1.20 (95\% CI 0.28--5.18, p=0.81). Subgroup analyses were performed to explore potential effect modification.

\subsection{Forest Plot: BRAF OR by Subgroup}
Forest plot available in the PDF report (Figure \ref{fig:braf_forest}). Key results:
\begin{itemize}
  \item All patients: OR=1.20 (p=0.81)
  \item N+ subgroup: OR=1.12 (p=0.89)
  \item N0 subgroup: not estimable (1 event)
  \item Multifocal: OR=1.45 (p=0.75)
  \item Unifocal: OR=1.32 (p=0.77)
  \item T1-T2: OR=1.02 (p=0.98)
  \item T3+: not estimable (1 event)
\end{itemize}

\subsection{Recurrence by BRAF and N Status}
\begin{table}[H]
\centering
\caption{Recurrence rates stratified by BRAF V600E and N status}
\label{tab:braf_n}
\begin{tabular}{lccc}
\toprule
Group & N & Recurrence & Rate \\
\midrule
BRAF$-$ / N0 & 38 & 1 & 2.6\% \\
BRAF$+$ / N0 & 17 & 0 & 0.0\% \\
BRAF$+$ / N$+$ & 28 & 3 & 10.7\% \\
BRAF$-$ / N$+$ & 31 & 3 & 9.7\% \\
\bottomrule
\end{tabular}
\end{table}
Recurrence rates are driven by N status, not BRAF. In N+ patients, BRAF+ (10.7\%) vs BRAF$-$ (9.7\%) are nearly identical.

\subsection{Mutation Co-occurrence}
\begin{table}[H]
\centering
\caption{Co-mutation counts among 114 patients}
\label{tab:comut}
\begin{tabular}{lccc}
\toprule
 & BRAF V600E & TERT & TP53 \\
\midrule
BRAF V600E & 45 & 2 & 0 \\
TERT & 2 & 6 & 1 \\
TP53 & 0 & 1 & 2 \\
\bottomrule
\end{tabular}
\end{table}
Very low co-occurrence: BRAF+TERT only 2 cases (1.8\%), BRAF+TP53: 0 cases.

\subsection{Interpretation}
BRAF V600E does not independently predict recurrence in this dataset. Possible explanations: (1) insufficient power (only 7 events), (2) BRAF effect may depend on histologic subtype (tall cell, hobnail) unavailable here, (3) BRAF may require TERT co-mutation for risk elevation (only 2 patients).

\section{Key Findings}
\begin{enumerate}
  \item \textbf{Firth penalized logistic regression} was applied to address complete separation. All variables now have finite odds ratios and confidence intervals.
  \item The overall recurrence rate was 6.1\% (7/114 patients).
  \item \textbf{N+ (lymph node metastasis)} was the strongest and most consistent predictor across all models: Firth OR=4.41 in univariate (p=0.112) and OR=4.98 in Model 1 (p=0.077).
  \item \textbf{TP53 mutation} (OR=2.81) showed a possible risk elevation but CI is extremely wide (0.06--124.82). \textbf{TERT} (OR=1.04) showed no association.
  \item \textbf{BRAF V600E} (prevalence 39.5\%) was not associated with recurrence in any subgroup (Firth OR=1.20; p=0.81 overall). Subgroup analyses (N$\pm$, multifocality, T stage) all show OR$\approx$1. Recurrence is driven by N+ status, not BRAF.
  \item All 4 multivariate models pass the Hosmer-Lemeshow goodness-of-fit test (all p$>$0.05).
  \item \textbf{Limitations:} Only 7 events. Firth resolves separation but cannot increase effective sample size. CIs for rare mutations remain very wide. Larger studies with Cox regression are recommended.
\end{enumerate}

\end{document}
